% ch22 — Color Science and Post-Processing
\chapter{Color Science and Post-Processing}
\label{ch:color-science}

The raw output of the ray tracer is a map of spectral intensities; turning
these into a displayable image requires a careful colour pipeline. This chapter
develops the physics of blackbody radiation, CIE colour matching, and the
tone-mapping and post-processing stages that produce the final render.

\section{Blackbody Radiation and the Planck Spectrum}
\label{sec:col-blackbody}
\coderef{Render.Blackbody}{hs}

We derive the Planck spectral radiance, compute the effective temperature from
the accretion-disk emissivity, and sample the spectrum at discrete wavelengths
for colour integration.

\section{CIE Colour Matching Functions}
\label{sec:col-cie}
\coderef{Render.Color}{hs}

We introduce the CIE 1931 standard observer, define the colour matching
functions $\bar{x}(\lambda)$, $\bar{y}(\lambda)$, $\bar{z}(\lambda)$, and
compute tristimulus values by integrating the spectral radiance against them.

\section{Spectral Integration}
\label{sec:col-spectral}

We describe the quadrature scheme that integrates the product of the Planck
function and the colour matching functions over the visible spectrum, balancing
accuracy against the number of wavelength samples.

\section{XYZ to sRGB Conversion}
\label{sec:col-xyz-srgb}

We apply the linear transformation from CIE XYZ to the sRGB colour space,
perform gamut clipping, and apply the sRGB gamma transfer function.

\section{ACES Tone Mapping}
\label{sec:col-tonemapping}

We implement the Academy Color Encoding System (ACES) tone-mapping curve that
compresses the high dynamic range of astrophysical intensities into the
displayable $[0,1]$ range while preserving perceptual contrast.

\section{Bloom, Glare, and Photon-Ring Smoothing}
\label{sec:col-bloom}

We describe the post-processing convolution filters that simulate optical
bloom and glare, and the targeted smoothing applied to photon-ring artefacts.

\section{GPU Shaders\starmark{}}
\label{sec:col-gpu-shaders}
\coderef{eh-gpu}{shaders}

We outline the fragment and compute shaders that perform tone mapping and
post-processing on the GPU, enabling real-time interactive exploration.

\paragraph{Key References.}
Blackbody physics follows \cite{Rybicki1979}; CIE colour science follows
\cite{Wyszecki1982}; ACES tone mapping follows the Academy specification
\cite{ACES2014}; bloom filters draw on \cite{Pharr2016}.
